% ===========================================================================
%  CHAPITRE 2 — Analyse et conception
% ===========================================================================
\thispagestyle{empty}
\begin{ChapterAbstract}
\end{ChapterAbstract}

\newpage

% ---------------------------------------------------------------------------
\section{Modélisation des acteurs et cas d'utilisation}

\subsection{Identification des acteurs}

Dans le cadre du Système Intégré de Gestion de Bibliothèque (SIGB) développé
pour l'EMP, plusieurs acteurs interagissent avec le système selon leurs rôles
et leurs responsabilités fonctionnelles (figure~\ref{fig:usercaseactors}).
L'identification de ces acteurs permet de structurer les besoins et de définir
clairement les interactions avec le système.

\begin{figure}[H]
    \centering
    \includegraphics[width=0.85\textwidth]{chapter_2/figures/usecase.png}
    \caption{Diagramme des cas d'utilisation — interactions avec le système}
    \label{fig:usercaseactors}
\end{figure}

Les principaux acteurs du système sont les suivants :

\begin{itemize}
    \item \textbf{Adhérent :} utilisateur final du système, il représente le
    lecteur inscrit à la bibliothèque. Il peut consulter ses informations,
    effectuer des réservations de documents et suivre l'état de ses emprunts.
\newpage
    \item \textbf{Agent de bibliothèque :} utilisateur opérationnel chargé de
    la gestion quotidienne des opérations de prêt. Il assure l'enregistrement
    des prêts, la restitution des documents, le renouvellement des emprunts
    ainsi que la consultation des fiches adhérents.

    \item \textbf{Administrateur :} acteur disposant de privilèges étendus sur
    le système. Il est responsable de la gestion des adhérents (ajout,
    modification, suppression), de la configuration des paramètres du système
    tels que les catégories d'adhérents, les pénalités et les postes.

    \item \textbf{Système planifié (tâche automatique) :} composant système
    exécuté de manière périodique via le planificateur de tâches Windows. Il
    assure le traitement automatique des pénalités, la mise à jour des états
    des adhérents ainsi que la gestion des réservations expirées.

    \item \textbf{Serveur d'authentification :} système externe à l'application
    assurant l'authentification centralisée des utilisateurs et la gestion des
    accès sécurisés au SIGB.
\end{itemize}

Ces acteurs constituent l'ensemble des entités interagissant avec le système
et permettent de couvrir l'ensemble des fonctionnalités métiers du SIGB,
allant des opérations de prêt jusqu'à la gestion administrative et automatisée.

% ---------------------------------------------------------------------------
\section{Modèle de données}

\subsection{Présentation de la base de données}

La base de données du SIGB de l'EMP est hébergée sous \textbf{Oracle Database
11g}. Elle a été conçue selon la méthode \textbf{MERISE} et comporte
\textbf{61 tables} réparties en cinq groupes fonctionnels distincts. Ce schéma
est le fruit d'un travail de conception entamé en 2007 : une partie des tables
est issue de la migration de l'ancien logiciel SYNGEB, l'autre partie ayant été
ajoutée pour répondre aux besoins spécifiques de l'EMP.

Le tableau~\ref{tab:groupes_tables} résume l'organisation thématique des
61 tables.
\newpage
\begin{table}[H]
\centering
\caption{Organisation thématique des 61 tables de la base de données CATLIB}
\label{tab:groupes_tables}
\renewcommand{\arraystretch}{1.4}
\begin{tabular}{|p{0.5cm}|p{4.5cm}|p{8cm}|}
\hline

\textbf{N°} & \textbf{Groupe} & \textbf{Tables principales} \\
\hline
1 & Notices bibliographiques &
\texttt{NOTICE}, \texttt{TYPE\_NOTICE}, \texttt{SOURCE\_ARTICLE},
\texttt{PERIODICITE}, \texttt{NOTICE\_COLLECTION}, \texttt{NEWACQUIS} \\
\hline
2 & Circulation et adhérents &
\texttt{ADHERENT}, \texttt{CATEGORIE}, \texttt{EXEMPLAIRE},
\texttt{ETAT\_EXEMPLAIRE}, \texttt{PRET}, \texttt{RESERVATION},
\texttt{HISTORIQUE\_PRET}, \texttt{PENALITE}, \texttt{PENALITE\_ADHERENT},
\texttt{PENALITE\_ADHERENT\_TEMP},
\texttt{HISTORIQUE\_PENALITE\_ADHERENT},
\texttt{JOURS\_FERIES},
\texttt{PARAMETRES\_CATLIB\_PRET},
\texttt{ETAT\_ADHERENT},
\texttt{POSITION}, \texttt{OPERATION} \\
\hline
3 & Indexation textuelle &
\texttt{TERME}, \texttt{TERME\_EXACT}, \texttt{NOTICE\_TERME},
\texttt{NOTICE\_TERME\_EXACT}, \texttt{MOTS\_CLES}, \texttt{MOTS\_VIDES},
\texttt{NOTICE\_MOT\_CLE}, \texttt{TABLE\_CDD}, \texttt{THEME},
\texttt{NOTICE\_THEME} \\
\hline
4 & Thèses et mémoires &
\texttt{NOTICE\_DIP\_DIS\_ETAB}, \texttt{DIPLOME},
\texttt{DISCIPLINE}, \texttt{ETABLISSEMENT} \\
\hline
5 & Informations bibliographiques &
\texttt{EDITEUR}, \texttt{VILLE}, \texttt{PAYS}, \texttt{PAYS\_PUBLICATION},
\texttt{LANGUE}, \texttt{NOTICE\_LANGUE}, \texttt{NOTICE\_EDITION},
\texttt{MENTION\_RESPONSABILITE}, \texttt{AUTEUR},
\texttt{AUTEUR\_SECONDAIRE}, \texttt{CO\_AUTEUR},
\texttt{MENTION\_RES\_COLLECTION}, \texttt{COLLECTION},
\texttt{SELECTION}, \texttt{SELECTION\_NOTICE} \\
\hline
\end{tabular}
\end{table}

\FloatBarrier

\subsection{Schéma des tables principales du module de prêt}

Les sections suivantes présentent le schéma conceptuel de la base de données
organisé par groupe fonctionnel, tel qu'il a été extrait depuis Oracle SQL
Developer. Chaque partie du schéma met en évidence les tables et leurs
relations dans un périmètre fonctionnel précis.

\subsubsection{Groupe 1 — Notices bibliographiques}

Le premier groupe rassemble les tables décrivant les ressources documentaires
gérées par la bibliothèque. La table centrale est \texttt{NOTICE}, qui contient
25 attributs couvrant toutes les métadonnées d'un document : identifiant unique
(\texttt{ID\_NOTICE}), type de notice (\texttt{ID\_TYPE}), titre propre,
sous-titre, cote de localisation, code CDD, ISBN/ISSN, résumé, nombre
d'exemplaires, accessibilité et type de données pour les ressources
électroniques. Elle entretient plus de 20 relations directes avec les autres
tables du schéma, et des relations indirectes via le champ \texttt{COTE}.
Les types de notices supportés sont : publication en série, mémoire/thèse,
monographie, article, tiré-à-part et ressource électronique.

\begin{figure}[H]
    \centering
    \includegraphics[width=\textwidth]{chapter_2/figures/notice.png}
    \caption{Schéma conceptuel — Partie 1 : notices bibliographiques}
    \label{fig:schema_p1}
\end{figure}

\FloatBarrier

\subsubsection{Groupe 2 — Circulation et adhérents}

Ce groupe constitue le cœur fonctionnel du module de prêt. Il regroupe
l'ensemble des tables qui permettent de gérer le cycle de vie d'un document
emprunté, depuis l'inscription de l'adhérent jusqu'à l'archivage de
l'historique des pénalités.
\newpage
\begin{itemize}
    \item \textbf{ADHERENT} : contient les informations de chaque lecteur
    (\texttt{ID\_ADHERENT}, \texttt{NOM}, \texttt{PRENOM},
    \texttt{ID\_POSITION}, \texttt{ID\_CATEGORIE}, \texttt{ETAT\_ADHERENT}).
    L'état encode le statut courant de l'adhérent (actif, pénalisé, suspendu).

    \item \textbf{CATEGORIE} : définit les paramètres de prêt par profil
    d'adhérent — nombre maximal de documents autorisés
    (\texttt{NOMBRE\_DOCUMENT}) et durée standard du prêt en jours
    (\texttt{DUREE\_PRET}).

    \item \textbf{EXEMPLAIRE} : représente un exemplaire physique d'un
    document (\texttt{ID\_EXEMPLAIRE}, \texttt{ID\_ETAT}, \texttt{COTE}).
    La table \texttt{ETAT\_EXEMPLAIRE} référence les états possibles
    (disponible, emprunté, réservé, exclu du prêt).

    \item \textbf{PRET} : enregistre chaque prêt actif avec
    \texttt{ID\_ADHERENT}, \texttt{ID\_EXEMPLAIRE}, \texttt{DATE\_PRET} et
    \texttt{ETAT\_DUREE} (indicateur de dépassement).

    \item \textbf{HISTORIQUE\_PRET} : archive les prêts clôturés avec la date
    effective de retour (\texttt{DATE\_RETOUR}).

    \item \textbf{RESERVATION} : gère les réservations actives.
    \texttt{HEURE\_RESERVATION} permet d'ordonner la file d'attente
    chronologiquement ; la cote du document reservé (\texttt{COTE}) sert de lien
    indirect avec \texttt{NOTICE}.

    \item \textbf{PENALITE} : barème des pénalités par catégorie d'adhérent
    et par palier de retard (\texttt{JOURS\_RETARD} →
    \texttt{NOMBRE\_JOURS\_RETARD}). Ce barème est progressif.

    \item \textbf{PENALITE\_ADHERENT} : pénalités actives appliquées à un
    adhérent (\texttt{DATE\_PENALITE}, \texttt{NOMBRE\_JOURS\_PENALITE}).

    \item \textbf{PENALITE\_ADHERENT\_TEMP} : table de travail temporaire
    utilisée par la tâche planifiée lors du calcul des pénalités.

    \item \textbf{HISTORIQUE\_PENALITE\_ADHERENT} : archive des pénalités
    expirées, permettant de conserver une traçabilité complète.

    \item \textbf{JOURS\_FERIES} : liste des jours fériés
    (\texttt{DATE\_JOUR\_FERIE}), utilisée lors du calcul des dates de retour
    prévues afin d'exclure ces jours du décompte.

    \item \textbf{PARAMETRES\_CATLIB\_PRET} : paramètre global de la durée
    de réservation (\texttt{DUREE\_RESERVATION}), c'est-à-dire le délai
    accordé à un adhérent pour venir récupérer un document réservé
    devenu disponible.

    \item \textbf{HISTORIQUE\_AUTH} : journal des opérations effectuées par
    les agents (\texttt{ID\_ADMIN}, \texttt{DATE\_OPERATION},
    \texttt{ID\_ADHERENT}, \texttt{ID\_TYPE\_OPERATION}), permettant un
    audit complet des actions sur le système.
\end{itemize}

\begin{figure}[H]
    \centering
    \includegraphics[width=\textwidth]{chapter_2/figures/circulation.png}
    \caption{Schéma conceptuel — Partie 2 : circulation et adhérents}
    \label{fig:schema_p2}
\end{figure}

\FloatBarrier

\subsubsection{Groupe 3 — Indexation textuelle}

Ce groupe supporte la recherche documentaire dans le portail OPAC. Il permet
d'associer à chaque notice des termes d'indexation (\texttt{TERME},
\texttt{TERME\_EXACT}), des mots-clés (\texttt{MOTS\_CLES}), des thèmes
(\texttt{THEME}), et des codes de la Classification Décimale de Dewey
(\texttt{TABLE\_CDD}). Les tables de liaison \texttt{NOTICE\_TERME},
\texttt{NOTICE\_TERME\_EXACT}, \texttt{NOTICE\_MOT\_CLE} et
\texttt{NOTICE\_THEME} assurent les associations many-to-many entre les notices
et ces entrées d'indexation. La table \texttt{MOTS\_VIDES} liste les mots
exclus de l'index de recherche.

\begin{figure}[H]
    \centering
    \includegraphics[width=0.9\textwidth]{chapter_2/figures/notice_terme.png}
    \caption{Schéma conceptuel — Partie 3 : indexation textuelle}
    \label{fig:pays}
\end{figure}

\FloatBarrier
\newpage
\subsubsection{Groupe 4 — Thèses et mémoires}

Ce groupe gère les métadonnées spécifiques aux notices de type thèse ou
mémoire, qui ne peuvent pas être entièrement décrites par la table
\texttt{NOTICE} seule. La table \texttt{NOTICE\_DIP\_DIS\_ETAB} joue le rôle
de table de liaison enrichie : elle rattache une notice à un diplôme
(\texttt{DIPLOME}), à une discipline scientifique (\texttt{DISCIPLINE}), à un
établissement de soutenance (\texttt{ETABLISSEMENT}), et conserve l'année de
soutenance (\texttt{ANNE\_SOUTENANCE}).

\begin{figure}[H]
    \centering
    \includegraphics[width=0.75\textwidth]{chapter_2/figures/diplome.png}
    \caption{Schéma conceptuel — Partie 4 : thèses et mémoires}
    \label{fig:schema_p4}
\end{figure}

\FloatBarrier

\subsubsection{Groupe 5 — Informations bibliographiques complémentaires}

Ce groupe décrit les entités bibliographiques d'accompagnement : pays et villes
de publication (\texttt{PAYS}, \texttt{VILLE}), éditeurs (\texttt{EDITEUR}),
langues (\texttt{LANGUE}), et adresses bibliographiques (\texttt{NOTICE\_EDITION}).
Il gère également les mentions de responsabilité des auteurs
(\texttt{MENTION\_RESPONSABILITE}), déclinées en auteur principal
(\texttt{AUTEUR}), co-auteurs (\texttt{CO\_AUTEUR}), auteurs secondaires
(\texttt{AUTEUR\_SECONDAIRE}) et mentions associées aux collections
(\texttt{MENTION\_RES\_COLLECTION}). La table \texttt{FONCTION} précise le
rôle d'un auteur secondaire (directeur, traducteur, illustrateur, etc.).

\begin{figure}[H]
    \centering
    \includegraphics[width=0.85\textwidth]{chapter_2/figures/pays.png}
    \caption{Schéma conceptuel — Partie 5a : publication (ville, pays, éditeur, langues)}
    \label{fig:schema_p5a}
\end{figure}

\begin{figure}[H]
    \centering
    \includegraphics[width=0.85\textwidth]{chapter_2/figures/auteur.png}
    \caption{Schéma conceptuel — Partie 5b : auteurs et mentions de responsabilité}
    \label{fig:schema_p5b}
\end{figure}

\FloatBarrier

% ---------------------------------------------------------------------------
\section{Opérations essentielles du module de prêt}

Le module de prêt constitue le cœur fonctionnel du système de gestion de
bibliothèque. Il orchestre l'ensemble des opérations liées à la circulation
des documents entre la bibliothèque et les adhérents, en respectant les règles
de gestion définies par l'établissement. Ces opérations reposent sur une série
de contrôles métier assurant la cohérence, la sécurité et la traçabilité des
transactions.

\subsection{Prêt d'un document}

Le processus de prêt débute par la vérification de l'adhérent, afin de
s'assurer de son éligibilité : état actif, absence de pénalité ou de
suspension, et respect du quota de documents autorisé selon sa catégorie
(champ \texttt{NOMBRE\_DOCUMENT} de la table \texttt{CATEGORIE}). Une fois
cette validation effectuée, le système procède à la vérification du document
demandé — en contrôlant son existence dans la base, la disponibilité d'un
exemplaire (\texttt{ETAT\_EXEMPLAIRE}), et l'éventuelle existence de
réservations ou d'emprunts en cours (figure~\ref{fig:diagpret}).

En fonction de ces vérifications, le système détermine l'action appropriée :
attribution directe d'un exemplaire disponible, affectation d'un exemplaire
réservé pour l'adhérent courant, proposition de réservation, ou refus de
l'opération. Le prêt est ensuite enregistré dans la table \texttt{PRET} avec
calcul automatique de la date de retour prévue en fonction de la durée de prêt
de la catégorie de l'adhérent (\texttt{DUREE\_PRET}) et des jours fériés
configurés dans la table \texttt{JOURS\_FERIES}.

\begin{figure}[H]
    \centering
    \includegraphics[width=0.85\textwidth]{chapter_2/figures/pret3.png}
    \caption{Diagramme d'activité — opération de prêt}
    \label{fig:diagpret}
\end{figure}

\FloatBarrier

\subsection{Restitution (retour)}

Le processus de restitution intervient lorsqu'un adhérent rapporte un document
emprunté. Cette opération vise non seulement à libérer l'exemplaire, mais aussi
à régulariser la situation de l'adhérent et à gérer les priorités de la file
d'attente des réservations. Le traitement suit plusieurs étapes clés
(figure~\ref{fig:diagrest}) :

\begin{itemize}
    \item \textbf{Validation initiale} : le système vérifie l'existence et la
    validité du prêt ainsi que l'identité de l'adhérent, pour s'assurer que
    les données correspondent aux enregistrements de la base (\texttt{PRET}).

    \item \textbf{Traitement des retards} : une vérification de la date
    d'échéance est effectuée. En cas de retard détecté, le système calcule
    automatiquement le nombre de jours de dépassement et applique les pénalités
    correspondantes au dossier de l'adhérent selon le barème de la table
    \texttt{PENALITE}, les enregistrant dans \texttt{PENALITE\_ADHERENT}.

    \item \textbf{Archivage du prêt} : le prêt est supprimé de la table
    \texttt{PRET} et archivé dans \texttt{HISTORIQUE\_PRET} avec la date de
    retour effective.

    \item \textbf{Mise à jour du statut du document} : le système analyse l'état
    des réservations en attente pour la cote concernée. Si une réservation est
    en attente, l'exemplaire passe au statut \og Réservé \fg{} et la
    notification de disponibilité est déclenchée. S'il n'existe aucune
    réservation, l'exemplaire repasse au statut \og Disponible \fg.
\end{itemize}

\begin{figure}[H]
    \centering
    \includegraphics[width=0.85\textwidth]{chapter_2/figures/rest2.png}
    \caption{Diagramme d'activité — opération de restitution}
    \label{fig:diagrest}
\end{figure}

\FloatBarrier

\subsection{Réservation}

La réservation permet à un adhérent de signaler son intention d'emprunter un
document actuellement indisponible (déjà prêté). Le système enregistre la
demande dans la table \texttt{RESERVATION} avec un horodatage
(\texttt{HEURE\_RESERVATION}) qui permet d'ordonner les demandes par ordre
d'arrivée et de gérer équitablement la file d'attente.

Lors de la restitution du document concerné, la réservation en tête de file est
traitée en priorité : l'exemplaire lui est affecté, son état est mis à
\og Réservé \fg{} et l'adhérent est notifié. Celui-ci dispose d'un délai
paramétrable (\texttt{DUREE\_RESERVATION} dans la table
\texttt{PARAMETRES\_CATLIB\_PRET}) pour se présenter et effectuer le prêt.
Passé ce délai, la tâche planifiée libère l'exemplaire et le propose au
réservataire suivant dans la file d'attente.

Un même document peut faire l'objet de réservations multiples, chacune liée à
un adhérent différent, formant ainsi une file gérée chronologiquement.

\FloatBarrier

\subsection{Renouvellement}

Le renouvellement permet à un adhérent de prolonger la durée d'un emprunt en
cours sans avoir à rapporter physiquement le document. Cette opération est
déclenchée soit depuis l'écran de restitution (bouton \og Renouvellement \fg{}),
soit directement depuis la fiche de l'adhérent.

Le système vérifie avant tout que le renouvellement est autorisé : l'adhérent
ne doit pas être pénalisé, et l'exemplaire ne doit pas faire l'objet d'une
réservation active par un autre adhérent. Si ces conditions sont réunies, le
prêt existant est clôturé et un nouveau prêt est immédiatement créé à la date
du jour, avec une nouvelle date de retour prévue calculée selon la durée de la
catégorie de l'adhérent et les jours fériés. L'opération est tracée dans
\texttt{HISTORIQUE\_AUTH}.

\FloatBarrier

\subsection{Gestion des pénalités}

Les pénalités constituent le mécanisme de régulation de la discipline de prêt.
Elles sont déclenchées automatiquement lorsqu'un adhérent dépasse la date de
retour prévue d'un document emprunté. Le système applique un barème progressif
défini dans la table \texttt{PENALITE}, qui associe à chaque catégorie
d'adhérent (\texttt{ID\_CATEGORIE}) un nombre de jours de pénalité
(\texttt{NOMBRE\_JOURS\_RETARD}) en fonction du nombre de jours de retard
(\texttt{JOURS\_RETARD}).

Lors de l'application d'une pénalité, l'enregistrement est créé dans la table
\texttt{PENALITE\_ADHERENT} avec la date de début et le nombre de jours de
sanction. L'état de l'adhérent (\texttt{ETAT\_ADHERENT}) est mis à jour pour
bloquer tout nouvel emprunt. La pénalité reste active jusqu'à son expiration,
calculée quotidiennement par la tâche planifiée. À l'expiration, l'enregistrement
est archivé dans \texttt{HISTORIQUE\_PENALITE\_ADHERENT} et l'adhérent est
automatiquement réactivé.

L'écran des paramètres de l'ancienne application (figure~\ref{fig:catlib_params})
illustre la saisie de ce barème : par exemple, pour la catégorie A, 1 jour de
retard entraîne 1 jour de pénalité, 30 jours de retard entraînent 90 jours,
et 60 jours de retard entraînent 365 jours de suspension.

\FloatBarrier

% ---------------------------------------------------------------------------
\section{La tâche planifiée Windows — programme console}
\label{sec:tache_planifiee}

\subsection{Rôle et objectifs}

Afin de garantir l'intégrité des règles de gestion sans intervention humaine
constante, le système intègre un \textbf{programme console .NET} exécuté
quotidiennement via le Planificateur de tâches Windows. Ce composant assure
deux fonctions essentielles que l'application web Blazor ne peut pas prendre
en charge de manière fiable, car elles ne sont pas liées à une session
utilisateur active :

\begin{itemize}
    \item \textbf{Mise à jour des pénalités} : détection quotidienne des prêts
    en retard, calcul et application des sanctions correspondantes selon le
    barème paramétré, et levée automatique des pénalités expirées avec
    réactivation de l'adhérent concerné.

    \item \textbf{Gestion des réservations expirées} : détection des réservations
    dont le délai d'attente (\texttt{DUREE\_RESERVATION}) est dépassé sans
    que l'adhérent soit venu emprunter le document, libération de l'exemplaire
    et réorganisation de la file d'attente au profit du réservataire suivant.
\end{itemize}

Ce découplage entre l'application interactive et le traitement automatisé
est une décision architecturale délibérée : le programme console est léger,
indépendant, facile à monitorer via les journaux Windows, et peut être exécuté
même lorsque aucun utilisateur n'est connecté au système Blazor.

\subsection{Logique de traitement quotidien}

\subsubsection*{Phase 1 — Analyse des prêts et détection des retards}

La première phase consiste en un parcours exhaustif de l'ensemble des prêts
enregistrés dans la table \texttt{PRET} afin d'identifier les dépassements de
délais et d'appliquer les règles de gestion associées. Pour chaque prêt, le
système effectue les opérations suivantes :

\begin{itemize}
    \item \textbf{Vérification de la date d'échéance} : la date de retour
    prévue est calculée à partir de \texttt{DATE\_PRET} et de
    \texttt{DUREE\_PRET} (catégorie de l'adhérent), en excluant les jours
    fériés de la table \texttt{JOURS\_FERIES}. Cette date est comparée à la
    date courante.

    \item \textbf{Détection des retards} : si la date actuelle dépasse la date
    de retour prévue, le prêt est marqué comme en retard (\texttt{ETAT\_DUREE}
    = \texttt{'R'}).

    \item \textbf{Application des pénalités} : le nombre de jours de
    dépassement est calculé et le barème correspondant est recherché dans la
    table \texttt{PENALITE}. La pénalité résultante est enregistrée dans
    \texttt{PENALITE\_ADHERENT} et l'état de l'adhérent est mis à jour dans
    \texttt{ETAT\_ADHERENT}.

    \item \textbf{Gestion des réservations expirées} : pour chaque exemplaire
    en état \og Réservé \fg, le système vérifie si le délai d'attente accordé
    au réservataire est dépassé. Si c'est le cas, la réservation est supprimée
    et la suivante dans la file est activée. Si aucune réservation suivante
    n'existe, l'exemplaire repasse à l'état \og Disponible \fg.
\end{itemize}
\newpage
\begin{figure}[H]
    \centering
    \includegraphics[width=0.85\textwidth]{chapter_2/figures/phase1.png}
    \caption{Diagramme d'activité : détection des retards et
    application des pénalités}
    \label{fig:phase1}
\end{figure}

\FloatBarrier

\subsubsection*{Phase 2 — Expiration des pénalités}

La seconde phase est dédiée à la gestion du cycle de vie des pénalités actives.
Elle permet de lever automatiquement les restrictions appliquées aux adhérents
une fois la durée de sanction écoulée. Pour chaque pénalité enregistrée dans
\texttt{PENALITE\_ADHERENT}, le système exécute les opérations suivantes :

\begin{itemize}
    \item \textbf{Calcul de la date d'expiration} : la date de fin de pénalité
    est déterminée en ajoutant \texttt{NOMBRE\_JOURS\_PENALITE} à
    \texttt{DATE\_PENALITE}.

    \item \textbf{Vérification d'expiration} : si la date courante est
    supérieure ou égale à la date de fin calculée, la pénalité est considérée
    comme expirée.

    \item \textbf{Réactivation de l'adhérent} : le champ \texttt{ETAT\_ADHERENT}
    est remis à la valeur correspondant à un état actif, autorisant à nouveau
    les emprunts.

    \item \textbf{Archivage} : l'enregistrement est déplacé vers la table
    \texttt{HISTORIQUE\_PENALITE\_ADHERENT} afin de conserver une traçabilité
    complète des sanctions passées.

    \item \textbf{Nettoyage} : la pénalité expirée est supprimée de la table
    \texttt{PENALITE\_ADHERENT} pour maintenir une base de données cohérente.
\end{itemize}

\begin{figure}[H]
    \centering
    \includegraphics[width=0.85\textwidth]{chapter_2/figures/phase2.png}
    \caption{Diagramme d'activité : expiration des pénalités et
    réactivation des adhérents}
    \label{fig:phase2}
\end{figure}

\FloatBarrier

\subsection{Intégration avec le planificateur de tâches Windows}

Le programme console est un exécutable autonome produit par la compilation
du projet \texttt{.NET\,10} en mode \emph{self-contained}. Son intégration
dans l'infrastructure de l'EMP repose sur le \textbf{Planificateur de tâches
Windows} (\emph{Task Scheduler}), disponible nativement sur les systèmes
Windows Server et Windows 10/11.

La configuration de la tâche prévoit les éléments suivants :

\begin{itemize}
    \item \textbf{Déclencheur} : exécution quotidienne à une heure fixe
    (par exemple 00h30), choisie pour correspondre à une période de faible
    activité sur le réseau intranet de l'EMP.

    \item \textbf{Action} : lancement de l'exécutable
    \texttt{BiblioPenaliteJob.exe} en fournissant la chaîne de connexion
    Oracle en paramètre sécurisé, ou en la lisant depuis un fichier de
    configuration chiffré.

    \item \textbf{Conditions} : la tâche est configurée pour s'exécuter
    même si aucun utilisateur n'est connecté au serveur, avec les privilèges
    d'un compte de service dédié disposant des droits d'accès à la base Oracle.

    \item \textbf{Journalisation} : chaque exécution produit un fichier de
    log daté (\texttt{log\_YYYYMMDD.txt}) consignant le nombre de prêts
    traités, les pénalités appliquées, les réservations libérées, ainsi que
    les éventuelles erreurs rencontrées.

    \item \textbf{Reprise sur erreur} : en cas d'échec, le planificateur est
    configuré pour relancer la tâche jusqu'à trois fois à intervalles de
    30 minutes, garantissant ainsi la continuité du traitement même en cas de
    redémarrage du serveur Oracle ou d'un problème réseau transitoire.
\end{itemize}



\FloatBarrier

% ---------------------------------------------------------------------------
\section*{Conclusion du chapitre}
\addcontentsline{toc}{section}{Conclusion du chapitre }

Ce deuxième chapitre a posé les bases analytiques et conceptuelles du projet.
L'identification des acteurs et la modélisation des cas d'utilisation ont permis
de délimiter précisément le périmètre fonctionnel du module de prêt. La
présentation détaillée de la base de données Oracle en cinq groupes
fonctionnels a mis en évidence la richesse et la cohérence du schéma existant,
sur lequel la nouvelle application s'appuie sans modification. La description
des cinq opérations métier essentielles — prêt, restitution, réservation,
renouvellement et gestion des pénalités — accompagnée de leurs diagrammes UML,
a formalisé les règles de gestion à implémenter. Enfin, la conception de la
tâche planifiée, traitée comme un composant architectural à part entière, a
montré comment automatiser les traitements périodiques critiques de manière
robuste et indépendante de l'interface web. Le chapitre suivant exposera les
choix technologiques et l'architecture logicielle adoptés pour concrétiser
cette conception.